\begin{abstract}
Modern healthcare informatics suffers from an intractable architectural trilemma: preserving stringent patient privacy under statutory mandates (HIPAA, GDPR), ensuring tamper-evident cross-institutional auditability, and achieving high-throughput semantic interoperability across heterogeneous healthcare stakeholders (acute care hospitals, diagnostic laboratories, ambulatory pharmacies, payors). Monolithic consortium blockchains encounter severe throughput degradation, global consensus bottlenecks, and privacy leakage, while centralized Health Information Exchanges (HIEs) introduce single points of failure, administrative lock-in, and opaque access logging. Furthermore, the immutability of distributed ledgers directly collides with the statutory ``Right to be Forgotten'' under GDPR Article 17. To overcome these fundamental barriers, this paper introduces \textbf{BlockchainForkTree}, a federated multi-chain architecture grounded in a directed acyclic tree topology of specialized blockchain forks coordinated by an independent metadata repository chain with on-chain shared governance. We formalize the \textbf{Shared Governance \& Fork-Tree Hypothesis}: \textit{a hierarchically partitioned fork-tree topology governed by an on-chain consortium consensus protocol provides provably superior audit integrity, cryptographic fault isolation, and resilient HL7 FHIR interoperability compared to monolithic distributed ledgers and centralized registries.} We present the formal graph-theoretic model of the fork-tree arborescence, define genealogical lineage block height inheritance, and formulate a hybrid cryptographic storage protocol combining off-chain AES-256-GCM vaults with on-chain SHA-256 hash anchors. To resolve inter-chain operational silos, we formulate an asymmetric request-fulfillment state machine for HL7 FHIR R4 resources and specify deterministic graph traversal algorithms ($\text{BFS-EHR}$ and $\text{DFS-EHR}$) with formal proofs of longitudinal reconstruction completeness and causal ordering. Security is fortified through a multi-tier role-based access control (RBAC) matrix and a sovereign patient consent engine. Comprehensive empirical validation across a seven-chain network confirms zero data leakage, sub-second cryptographic verification, linear throughput scalability, and complete fault isolation under simulated Byzantine node failures.
\end{abstract}

\begin{IEEEkeywords}
Blockchain, Fork Tree Topology, HL7 FHIR, Healthcare Interoperability, Consortium Governance, Cryptographic Auditing, Off-chain Storage, HIPAA/GDPR Compliance, Multi-Chain Systems.
\end{IEEEkeywords}
